\ifdefined\standalone 
\documentclass[12pt]{article}
\usepackage{graphicx}
\usepackage{float}
\usepackage[a4paper, margin=1in]{geometry}
\usepackage[utf8]{inputenc}
\usepackage{textcomp}
\usepackage{amsmath}
\usepackage{booktabs}
\usepackage{tikz}
\usepackage{enumitem}
\usetikzlibrary{decorations.pathmorphing}
\usepackage[hidelinks]{hyperref} % <- hypertext

\begin{document}
\fi


\section{Extension to 2D and 3D Configurations}

To further evaluate the robustness and consistency of the solver, the reference one-dimensional cone calorimeter configuration is extended to two- and three-dimensional domains. The domain is expanded to $10~\mathrm{cm} \times 10~\mathrm{cm}$ in the $x$ and $y$ directions, while maintaining a vertical thickness of $18~\mathrm{mm}$ identical to the 1D case.

The computational mesh is composed of 20 uniform cells in both the $x$ and $y$ directions, and 180 cells in the $z$ direction. This corresponds to a spatial resolution of $\Delta x = \Delta y = 5~\mathrm{mm}$ and $\Delta z = 0.1~\mathrm{mm}$, which is 50 times finer along the vertical axis. A constant time step of $\Delta t = 0.5~\mathrm{s}$ is employed throughout the simulation.

To avoid the stability constraints typically imposed by the fine vertical resolution, the numerical scheme is configured to be implicit along the $z$-axis and explicit in the $x$ and $y$ directions. This is achieved by activating the option \texttt{SOLVER\_MAIN\_DIRECTION='Z'}.

In order to ensure theoretical equivalence with the one-dimensional configuration, adiabatic boundary conditions are imposed on all lateral surfaces (in the $x$ and $y$ directions). Under these assumptions, the 2D and 3D configurations are expected to reproduce the same thermal and chemical behavior as the 1D case.

The comparison of the Mass Loss Rate (MLR) obtained from the 1D, 2D, and 3D simulations is presented in Figure~\ref{fig:mlr_2d_3d}. The consistency between the results confirms the proper extension of the solver to multidimensional domains without introducing numerical bias.

\begin{figure}[H]
\centering
\includegraphics[width=0.75\textwidth]{mlr_2d_3d.png}
\caption{Comparison of Mass Loss Rate (MLR) between 1D, 2D, and 3D configurations}
\label{fig:mlr_2d_3d}
\end{figure}




\ifdefined\standalone
\end{document}
\fi






